\subsubsection{\textit{Exports service}}
\label{sec:exports-service}

El \textit{exports service} es el encargado de gestionar las versiones exportadas de las aplicaciones del ecosistema.
Concretamente, administra los archivos comprimidos en formato ZIP correspondientes a las dos aplicaciones del sistema:
\texttt{ftr\_world\_editor} y \texttt{ftr\_game}, para los sistemas operativos soportados, \textit{Linux} y
\textit{Windows}. Cada versión se identifica mediante un número de versión con formato semántico \texttt{vX.Y.Z} y
puede ser marcada como la versión más reciente (\textit{latest}) para su combinación de aplicación y sistema operativo.
\vspace{0.25cm}

La subida de una nueva versión se realiza mediante una petición \textit{multipart/form-data} a
\texttt{PUT /exports/zip}. La petición debe incluir los campos \texttt{app}, \texttt{version} y \texttt{os}, junto
con el archivo comprimido bajo el campo \texttt{file}. El sistema valida que todos estos campos y que el formato del archivo sea
efectivamente un ZIP. Si ya existe una entrada para la misma combinación de aplicación, versión y sistema operativo,
el sistema retorna un error de conflicto. De lo contrario, el archivo se almacena y se persiste el \textit{path}
resultante en la base de datos.
\vspace{0.25cm}

El sistema ofrece dos mecanismos para consultar versiones existentes. Por un lado, \texttt{GET /exports/zip} permite
obtener el \textit{path} de una versión específica indicando los parámetros \texttt{app}, \texttt{version} y
\texttt{os}; el parámetro \texttt{version} es opcional, en cuyo caso el sistema retorna el \textit{path} de la
versión marcada como \textit{latest} para esa combinación. Por otro lado, \texttt{GET /exports/zip/versions} retorna el 
listado completo de versiones disponibles, pudiendo filtrarse opcionalmente por \texttt{app} y \texttt{os}.
\vspace{0.25cm}

Para marcar una versión como la más reciente, se utiliza \texttt{PATCH /exports/zip/latest}, enviando en el cuerpo
de la petición los campos \texttt{app}, \texttt{version} y \texttt{os}. El sistema valida que la versión indicada
exista previamente y, de ser así, la marca como \textit{latest} para esa combinación de aplicación y sistema
operativo. Finalmente, una versión puede eliminarse mediante \texttt{DELETE /exports/zip}, indicando los parámetros 
\texttt{app}, \texttt{version} y \texttt{os}, los tres obligatorios.

Para más información donde se utilizan específicamente estos \textit{endpoints}, ver la sección de \nameref{sec:backoffice}.

\vspace{0.25cm}
